\begin{center}
    \large
    \begin{singlespace}    
        \textbf{\chineseTitle{}} \\[0.5cm]
    \end{singlespace}

    \begin{singlespace}    
    學生：\studentCnName  \hspace{2.5cm}  指導教授：\advisorCnName \hspace{0.1cm} 博士 \\

    \ifdefined\advisorCnNameB % 如果有共同指導教授
        \hspace{9.6cm}  \advisorCnNameB ~ 博士 \\
    \fi
    \end{singlespace}

    \vspace{0.5cm}

    \begin{singlespace}
        國立陽明交通大學\\
        \NameofDepartmentInstituteCN\\[0.2cm]
    \end{singlespace}

    \textbf{摘\hspace{1cm}要} \\[0.5cm]

\end{center}

\normalsize 

本研究針對台灣股票市場橫斷面報酬預測問題，以消融實驗框架系統性比較非圖基準模型、圖注意力網路（Graph Attention Network, GAT）及其中性化延伸模型，量化圖拓樸設計與訓練視窗長度對截面排序訊號品質與投資組合表現之影響。

實證資料採用台灣經濟新報（TEJ）台股日頻資料，以 56 維技術與估值特徵為模型輸入，預測目標為未來五日報酬之截面去均值。本文建構同產業完全連結之產業圖與全截面完全連結之全市場圖，並比較線性迴歸、XGBoost、長短期記憶網路（LSTM）、深度神經網路（DNN）及五種圖注意力變體共九類模型。三種訓練視窗（短期、中期、長期）採時序切分，以日頻資訊係數（Daily IC）、IC 資訊比率（ICIR）及投資組合績效指標進行評估；策略層面每五個交易日依預測分數選取前 10\% 股票等權做多，並以台灣 50 ETF（0050）為市場基準。

研究結果顯示，訓練視窗長度為樣本外可用性之主要決定因素：短期視窗下多數模型雖具正向日頻 IC，但投資組合年化報酬多接近零或為負；中、長期視窗下各模型 Sharpe 比率明顯改善。圖拓樸方面，GAT-industry 在日頻 IC 與 Sharpe 比率上均優於 GAT-universe，反映具有經濟意義之局部產業圖比全截面完全連結更能保留個股截面差異。GAT-IndNeutral 於長期視窗取得最高日頻 IC（0.0590）與 ICIR（0.786），GAT-full 取得最高 Sharpe 比率（1.849），兩者呈現截面排序訊號穩健性與投資組合絕對報酬之取捨。XGBoost 雖於長期視窗達最高總報酬（143.08\%），但其訓練集與測試集資訊係數差距懸殊（0.7824 對 0.0205），反映顯著之過擬合風險。整體而言，深度神經網路已展現特徵工程本身的排序效益，圖注意力機制之增量貢獻主要體現於長期視窗下訊號穩定性與拓樸可解釋性的提升。

\vspace{1cm}

% 中文摘要及關鍵詞 5-7 個
\textbf{關鍵字：}圖注意力網路、截面報酬預測、資訊係數、圖拓樸設計、台灣股票市場
